\documentclass[a4paper,11pt]{jsarticle}

\usepackage{amsmath,amssymb}
\usepackage{graphicx}
\usepackage{siunitx}
\usepackage{physics}
\usepackage{bm}
\usepackage{hyperref}

\title{電気回路Ⅲ 学習資料}
\author{}
\date{}

\begin{document}

\maketitle

\section{この資料の目的}

この資料は，電気回路Ⅲの内容を基礎から復習し，期末試験に向けて理解を整理するためのものである。

\section{前提知識}

ここでは，電圧，電流，抵抗，インピーダンス，フェーザ表示など，電気回路Ⅲを理解するために必要な基礎事項を整理する。

\section{過渡現象}

過渡現象では，回路にスイッチ操作や入力変化があったとき，電圧や電流が時間とともにどのように変化するかを扱う。

\section{二端子対回路}

二端子対回路では，入力端子と出力端子をもつ回路を，行列やパラメータを用いて表す。

\section{非正弦波交流}

非正弦波交流では，正弦波ではない周期波形をフーリエ級数などによって解析する。

\section{試験対策メモ}

出題範囲，重要公式，解法パターン，間違えやすい点を整理する。

\end{document}
